% Load the class with the 'code' option to typeset code snippets
% Then you must have minted installed and compile with the -shell-escape option
\documentclass[]{ccdsexam}

\usepackage{lipsum}

\begin{document}

\course{SC4012/CE4067/CZ4067}{SOFTWARE SECURITY}
\quiztitle{SEMESTER 2 EXAMINATION 2024-2025}
\monthyear{Apr/May 2025}
\length{2 hours}

%%%%%%%%%%%%%%%%%%%%%%%%
% Override exam instructions below if necessary
%\makeatletter
%\renewcommand{\printInstructions}{
%  \noindent{\bfseries\underline{INSTRUCTIONS}}
%
%  \begin{enumerate}[leftmargin=*,labelsep=20pt,topsep=0pt]
%    \item This paper contains \numquestions~questions and comprises \pageref{LastPage} pages.
%    \item Answer \textbf{ALL} questions.
%    \item This is a closed-book examination.
%  \end{enumerate}
%}
%\makeatother
%%%%%%%%%%%%%%%%%%%%%%%%

\makeheading

%%%%%%%%%%%%%%%%%%%%%%%%
% Do not print solutions by default
\noprintanswers
% To print solutions, use this instead
% \printanswers
%%%%%%%%%%%%%%%%%%%%%%%%

%%%%%%%%%%%%%%%%%%%%%%%%
% BEGIN EXAM QUESTIONS %
%%%%%%%%%%%%%%%%%%%%%%%%

\begin{questions}

%%%%%%%%%%%%%%%%%%%%%%%%%%%%%%%%% QUESTION %%%%%%%%%%%%%%%%%%%%%%%%%%%%%%%%%

\question \lipsum[1][1]

\begin{parts}

  \part[10] \label{part:1a} \lipsum[1][1-2]
  \droppoints

  \begin{solution}
    \lipsum[2][1-4]
  \end{solution}

  \part[15] \lipsum[1][3] As given in \cref{part:1a}, \lipsum[1][4-5]
  \droppoints

  \begin{solution}
    \lipsum[2][1-3]
  \end{solution}

\end{parts}

%%%%%%%%%%%%%%%%%%%%%%%%%%%%%%%%% QUESTION %%%%%%%%%%%%%%%%%%%%%%%%%%%%%%%%%

\question[25] \lipsum[1][1-3]

\begin{itemize}
  \item \lipsum[2][1]
  \item \lipsum[2][2]
  \item \lipsum[2][3]
\end{itemize}

\lipsum[1][4-6]
\droppoints

\begin{figure}[ht]
%%%%%%% A sample code snippet
%
%  \begin{minted}[breaklines]{c}
%    void dummy(int a)
%    {
%      int x = 1, y = 0;
%      assert(x - a != 0);
%    }
%  \end{minted}
%%%%%%%%%%%%%%%%%%%%%%%%%%%%%%%
  \centering
  \includegraphics[width=.5\textwidth]{dummy}
  \caption{A dummy figure.}\label{fig:test}
\end{figure}

\begin{figure}[ht]
  \centering
  \includegraphics[width=.5\textwidth]{dummy}
  \caption{}
\end{figure}

\begin{solution}
  \lipsum[2][1-4]
\end{solution}

% This marks the end of a question in case the footnote warning does not work correctly
\addquestionobject

%%%%%%%%%%%%%%%%%%%%%%%%%%%%%%%%% QUESTION %%%%%%%%%%%%%%%%%%%%%%%%%%%%%%%%%

\question

\begin{parts}
  \part \lipsum[1][1-3]
    \begin{subparts}
      \subpart[5] \lipsum[1][1] \label{part:format-read}
      \droppoints

      \begin{solution}
        \lipsum[2][1-2]
      \end{solution}

      \subpart[5] \lipsum[1][3]
      \droppoints

      \begin{solution}
        \lipsum[2][3-5]
      \end{solution}
    \end{subparts}

\part[15] \Cref{fig:dummy} \lipsum[1][4]

% Override the figure number to Q3 (instead of Q3a)
\renewcommand{\thefigure}{Q3}
\begin{figure}[ht]
  \centering
  \includegraphics[width=.5\textwidth]{dummy}
  \caption{}\label{fig:dummy}
\end{figure}

\begin{solution}
  \lipsum[2][3-5]
\end{solution}

\droppoints
% This marks the end of a question in case the footnote warning does not work correctly
\addquestionobject

\end{parts}

%%%%%%%%%%%%%%%%%%%%%%%%%%%%%%%%% QUESTION %%%%%%%%%%%%%%%%%%%%%%%%%%%%%%%%%

\question
\begin{parts}
  \part \lipsum[1][1-3]
  \begin{subparts}
  \subpart[8] \lipsum[1][1]
  \droppoints

  \begin{solution}
    \lipsum[2][1-2]
  \end{solution}

  \subpart[7] \lipsum[1][3]
  \droppoints

  \begin{table}[ht]
    \caption{}
    \centering
    \begin{tabular}{|c|c|c|c|c|}
      \hline
      \phantom{text} & \phantom{text} & \phantom{text} & \phantom{text} & \phantom{text} \\
      \hline
      \phantom{text} & \phantom{text} & \phantom{text} & \phantom{text} & \phantom{text} \\
      \hline
    \end{tabular}
  \end{table}

  \begin{solution}
    \lipsum[2][3-5]
  \end{solution}
  \end{subparts}

  \part[10] \lipsum[1][1-3]
  \droppoints

  \begin{solution}
    \lipsum[2][1-2]
  \end{solution}
\end{parts}

\end{questions}

\end{document}
